\section{Modèles représentatifs}

Hypothèses communes~:
\begin{itemize}
  \item Pas de panne
  \item Calculs locaux gratuits
  \item identifiants pour les sommets $u \mapsto \text{id}(u)$ sur $k$ bits où $k \geq \log_2(n)$ où $n$ est le nombre de sommets du graphe d'infrastructure.
\end{itemize}

\subsection{Modèle LOCAL}
\begin{itemize}
  \item hypothèses communes
  \item Communication synchrone
  \item taille de messages illimitée
\end{itemize}
Ce modèle permet d'étudier la nature locale d'un problème.

\subsection{Modèle CONGEST}
\begin{itemize}
  \item hypothèses communes
  \item Communication synchrone
  \item taille de messages bornée à $O(\log n)$ bits 
\end{itemize}

On peut aussi borner à la taille à $b \in \N$ bits. On parlerait alors de $b$-CONGEST.


\begin{exemple}
  Supposons vouloir détecter si deux voisins d'une machines sont voisins.
  \begin{lstLNat}
    Code de u
    Ronde 1 : SENDALL(id(u))
              L = liste des identifiants reçus par les voisins
    Ronde 2 : SENDALL(L)
              L' = union de toutes les L reçues par les voisins.
              Si id(u) est dans L', alors deux voisins de u sont voisins.
  \end{lstLNat}
  Cet algorithme, bien que valide, ne fonctionne pas dans le modèle CONGEST car la taille de L peut être de $O(n \log n)$ bits.
\end{exemple}

